\chapter{Considerações parciais e próximos passos}
\label{cap:consideracoes}

As Fases 1 e 2 transformaram percepções sobre o SIGAA/UnB em uma especificação de avaliação verificável. A Fase 1 delimitou o propósito, as partes interessadas e as decisões apoiadas, adaptou o modelo da ISO/IEC 25010:2023 ao contexto de uma avaliação caixa-preta e selecionou, por uma matriz Impacto~$\times$~Risco, quatro características: Capacidade de Interação, Adequação Funcional, Eficiência de Desempenho e Flexibilidade. A Fase 2 desdobrou essas características em quatro objetivos, 14 questões com hipóteses e 20 métricas, com níveis de pontuação e critérios de julgamento fixados antes da coleta, de modo que o resultado não possa ser ajustado aos dados.

As principais limitações já conhecidas são a amostra de conveniência, concentrada na FCTE; a medição de desempenho apenas do lado do cliente, sujeita a variações de rede e horário; a impossibilidade de avaliar matrícula e trancamento sem alterar dados reais; e a ausência de acesso ao código-fonte, que impede conclusões sobre a qualidade interna do produto.

A Entrega 2 cobrirá as Fases 3 e 4. Na Fase 3, a equipe produzirá o plano de avaliação: roteiro do moderador e termo de consentimento do teste com usuários, publicação do questionário, protocolo das medições com Lighthouse e das inspeções, ambiente de execução (dispositivos, navegadores e rede) e cronograma. Na Fase 4, as medidas serão obtidas, comparadas aos níveis de pontuação e julgadas; as hipóteses serão classificadas e as relações entre as características, discutidas. A tarefa e a característica de pior resultado orientarão a ação de melhoria obrigatória, como um protótipo de alta fidelidade para o fluxo pior avaliado, caso o problema seja de interface ou navegação. Na gestão do projeto, os marcos e as \textit{issues} serão criados para permitir o \textit{burndown}, e a planilha de esforço passará a alimentar os indicadores de recursos e a tabela de contribuição.
